\chapter{\PETSc}
\label{section:petsc}


% \begin{figure}[htb]
 \begin{center}
  \includegraphics[width=0.2\textwidth]{61_pidt/ExCALIBUR.png}
 \end{center}


\noindent
\Peano\ offers a \PETSc\ integration.
We mainly use \PETSc\ as prototyping backend for our studies around multigrid
ingredients.
These are funded by the UK's ExCALIBUR project.


\section{Configure \Peano\ with \PETSc}

\marginpar{todo: add info here how to configure system properly}

\section{Case study: Poisson equation solver on regular grid}


\begin{remark}
 This case study can be run via
 \begin{bash}
 cd myPeanoDir/benchmarks/multigrid/petsc/poisson
 export PYTHONPATH=../../../../python/
 python3 poisson.py
 \end{bash}
\end{remark}
In this initial case study, we solve the all-time classic in any introductory
Finite Element/Finite Differences book: 
The Poisson on the unit square ($d=2$) or cube ($d=3$) with the right-hand side

\begin{eqnarray*}
  - \Delta u & = & d \cdot \pi ^2 \prod _d \sin (\pi x ) \\
  u | _{\partial \Omega } & = & 0,
\end{eqnarray*} 

\noindent
for which we know the trivial analytic solution $u(x) = \prod _d \sin (\pi x )
$.
For this, we create a simple \PETSc\ project:

\begin{code}

\end{code}


\noindent
The project will yield the following trivial steps:
\begin{enumerate}
  \item Create grid (\Peano\ mesh sweep). This step might be repeated multiple
  times, as we only can add one grid level per grid sweep, and as we might also decide to load balance.
  \item Enumerate (\Peano\ mesh sweep). This step computes a local enumeration
  of the vertices (aka degrees of freedom).
  Eventually, it reduces the total degree of freedom count which allows us to
  map the local enumeration onto a global scheme.
  \item Set up \PETSc\ (\PETSc\ call).
  \item Init (\Peano\ mesh sweep). This step sets up the right-hand side of all
  mesh vertices. 
  \item Assemble (\Peano\ mesh sweep). This step actually sets up the equation
  system and creates the right-hand side, i.e.~it sets $A$ and $b$ of $Au=b$.
  \item Solve (\PETSc\ call).
  \item MapSolutionOntoMesh (\Peano\ mesh sweep). This takes the solution of $u$
  and writes is back into the vertices.
  \item Plot (\Peano\ mesh sweep). Plot the mesh plus its $u$-values.
\end{enumerate}




